% ===== CHƯƠNG 4: THỰC HIỆN VÀ KẾT QUẢ =====

\chapter{Thực Hiện và Kết Quả}

\section{Công nghệ sử dụng}

\subsection{Backend}

\begin{itemize}
  \item \textbf{Ngôn ngữ:} Python 3.10+ (sử dụng type hints, dataclass, walrus operator).
  \item \textbf{Thư viện:} Chỉ dùng Python Standard Library:
    \begin{itemize}
      \item \texttt{http.server} --- HTTP server threading.
      \item \texttt{urllib.parse} --- URL parsing.
      \item \texttt{json} --- JSON serialization.
      \item \texttt{math} --- Tính toán Haversine.
      \item \texttt{heapq} --- Priority queue cho thuật toán.
      \item \texttt{dataclasses} --- Station, Edge.
      \item \texttt{pathlib} --- Path handling.
    \end{itemize}
  \item \textbf{Database:} Không dùng. Dữ liệu lưu trong JSON file và load vào memory.
\end{itemize}

\subsection{Frontend}

\begin{itemize}
  \item \textbf{HTML5:} Semantic markup.
  \item \textbf{CSS3:} Custom properties, CSS Grid, Flexbox, media queries (responsive).
  \item \textbf{JavaScript:} ES2020 (const/let, arrow functions, template literals, spread operator, async/await).
  \item \textbf{Leaflet.js v1.9.4:} Interactive web map, markers, polylines, Popup.
  \item \textbf{OpenStreetMap:} Tile layer (không cần API key).
\end{itemize}

\subsection{Deployment \& Testing}

\begin{itemize}
  \item \textbf{Testing Framework:} Python \texttt{unittest} (built-in).
  \item \textbf{Deployment:} Standalone Python HTTP server (không cần app server như Gunicorn, nếu volume tăng có thể thêm).
  \item \textbf{Containerization:} (Optional) Docker.
\end{itemize}

\section{Cấu trúc mã nguồn}

\begin{lstlisting}[caption=Cấu trúc thư mục project]
Project-for-Introduction-to-AI/
|-- app.py                          # Entry point
|-- server.py                       # HTTP server + endpoints
|-- metro_app/
|   |-- __init__.py
|   |-- data.py                     # Station, Edge, load_network()
|   |-- algorithms.py               # UCS, Greedy, A*, Dijkstra
|   |-- service.py                  # Business logic (find_routes, etc.)
|   |-- geo.py                      # haversine_km()
|   `-- osm_import.py               # (Optional) Refresh data from Overpass API
|-- web/
|   |-- index.html
|   |-- app.js
|   `-- styles.css
|-- data/
|   `-- shanghai_metro_osm.json     # 411 ga, 20 tuyen (206 KB)
|-- tests/
|   |-- test_algorithms.py
|   `-- test_http_server.py
`-- README.md
\end{lstlisting}

\section{Thực hiện 4 thuật toán}

\subsection{Uniform Cost Search}

\begin{lstlisting}[language=Python, caption=UCS implementation (metro\_app/algorithms.py - excerpt)]
def uniform_cost_search(
    graph: dict[str, list[Edge]],
    start: str,
    goal: str,
) -> SearchResult:
    queue: list[tuple[float, int, str]] = [(0.0, 0, start)]
    tracker = count(1)
    came_from: dict[str, tuple[str, str] | None] = {start: None}
    best_cost = {start: 0.0}
    explored = 0

    while queue:
        current_cost, _, current = heapq.heappop(queue)
        if current_cost > best_cost[current]:
            continue

        explored += 1
        if current == goal:
            path, lines = _reconstruct_path(came_from, goal)
            return SearchResult("UCS", path, round(current_cost, 2), explored,
                                lines, _path_distance(graph, path, lines))

        for edge in graph[current]:
            new_cost = current_cost + edge.travel_time
            if edge.target not in best_cost or new_cost < best_cost[edge.target]:
                best_cost[edge.target] = new_cost
                came_from[edge.target] = (current, edge.line)
                heapq.heappush(queue, (new_cost, next(tracker), edge.target))

    raise ValueError(f"Cannot reach {goal} from {start}.")
\end{lstlisting}

\subsection{Dijkstra}

Dijkstra được cài đặt tương tự UCS nhưng tính toàn bộ SSSP trước khi reconstruct path, dẫn tới số node explored lớn hơn.

\subsection{Greedy Best-First Search}

\begin{lstlisting}[language=Python, caption=Greedy BFS pseudocode (thực hiện tương tự)]
def greedy_best_first_search(
    graph: dict[str, list[Edge]],
    stations: dict[str, Station],
    start: str,
    goal: str,
) -> SearchResult:
    # Dùng heuristic h(n) làm priority
    # Early-exit khi pop goal từ queue
    pass
\end{lstlisting}

\subsection{A* Search}

\begin{lstlisting}[language=Python, caption=A* pseudocode]
def a_star_search(
    graph: dict[str, list[Edge]],
    stations: dict[str, Station],
    start: str,
    goal: str,
) -> SearchResult:
    # Dùng f(n) = g(n) + h(n) làm priority
    # g(n) = chi phí từ start, h(n) = heuristic
    pass
\end{lstlisting}

\section{Hàm Heuristic}

\begin{lstlisting}[language=Python, caption=Hàm heuristic (metro\_app/algorithms.py)]
def heuristic(stations: dict[str, Station], current: str, goal: str) -> float:
    here = stations[current]
    destination = stations[goal]
    return haversine_km(here.lat, here.lon, destination.lat, destination.lon) / 35.0 * 60.0
\end{lstlisting}

Công thức: Khoảng cách Haversine (km) ÷ 35 (km/h) × 60 (phút).

\begin{lstlisting}[language=Python, caption=Hàm Haversine (metro\_app/geo.py)]
def haversine_km(lat1: float, lon1: float, lat2: float, lon2: float) -> float:
    from math import radians, sin, cos, sqrt, atan2

    R = 6371  # Earth radius in km

    phi1, phi2 = radians(lat1), radians(lat2)
    delta_phi = radians(lat2 - lat1)
    delta_lambda = radians(lon2 - lon1)

    a = sin(delta_phi/2)**2 + cos(phi1) * cos(phi2) * sin(delta_lambda/2)**2
    c = 2 * atan2(sqrt(a), sqrt(1-a))

    return R * c
\end{lstlisting}

\section{Kiểm thử}

Bộ test gồm 29 test case:

\subsection{Unit Tests (metro\_app tests)}

\begin{itemize}
  \item \textbf{14 tests thuật toán:} Mỗi thuật toán được test trên các đồ thị sample khác nhau (simple path, multiple paths, branching).
  \item \textbf{3 tests correctness:} Verify kết quả path, cost, explored count so sánh với expected.
  \item \textbf{2 tests Dijkstra:} So sánh SSSP đầy đủ vs goal-directed behavior.
  \item \textbf{Test filter:} Verify graph filter khi loại tuyến/đoạn.
  \item \textbf{Test geo bounds:} Verify nearest-station reject điểm > 50 km.
\end{itemize}

\subsection{Integration Tests (HTTP server tests)}

\begin{itemize}
  \item \textbf{4 endpoints × (happy path + malformed input):} Test \texttt{/api/network}, \texttt{/api/routes}, \texttt{/api/nearest-station}, \texttt{/api/routes-by-points}.
  \item \textbf{10 tests HTTP:} Gồm status code 200, 400, content-type JSON.
\end{itemize}

\begin{lstlisting}[language=bash, caption=Chạy kiểm thử]
$ python -m unittest discover -s tests -v
test_algorithms.py ............. [19 tests]
test_http_server.py ............ [10 tests]
Ran 29 tests in 0.6s
OK
\end{lstlisting}

\section{Kết quả Benchmark}

Chạy benchmark trên một cặp điểm: **Hongqiao Railway Station → Century Avenue**

% ===== PLACEHOLDER BẢNG =====
% TODO: Điền dữ liệu benchmark thực tế từ project
% Bảng so sánh cost, distance, explored nodes, path length
\begin{table}[H]
  \centering
  \caption{Kết quả benchmark so sánh 4 thuật toán}
  \label{tab:benchmark}
  \begin{tabular}{lrrrr}
    \toprule
    \textbf{Thuật toán} & \textbf{Cost (phút)} & \textbf{Distance (km)} & \textbf{Explored} & \textbf{Path Length} \\
    \midrule
    UCS & 48.61 & 21.82 & 170 & 9 ga \\
    Greedy & 48.85 & 21.96 & 15 & 9 ga \\
    A* & 48.61 & 21.82 & 64 & 9 ga \\
    Dijkstra & 48.61 & 21.82 & $\geq 170$ & 9 ga \\
    \bottomrule
  \end{tabular}
\end{table}

\subsection{Phân tích kết quả}

\begin{itemize}
  \item \textbf{UCS, A*, Dijkstra:} Tất cả đều tối ưu, chi phí bằng nhau (48.61 phút, 21.82 km).

  \item \textbf{Greedy:} Chi phí cao hơn (48.85 phút, 21.96 km) vì chỉ dùng heuristic mà không dùng chi phí thực tế.

  \item \textbf{Explored nodes:}
    \begin{itemize}
      \item Greedy: 15 (nhanh nhất, nhưng không tối ưu).
      \item A*: 64 (cân bằng tốt, tối ưu).
      \item UCS: 170 (không dùng heuristic, nhưng tối ưu).
      \item Dijkstra: $\geq 170$ (tính toàn bộ SSSP).
    \end{itemize}

  \item \textbf{Path length:} Bốn thuật toán đều tìm path 9 ga (do tính tối ưu hoặc sự trùng lặp đường đi).
\end{itemize}

% ===== PLACEHOLDER HÌNH =====
% TODO: Vẽ biểu đồ bar chart so sánh explored\_nodes của 4 algo
% Gợi ý file: figures/benchmark-chart-explored-nodes.png
% Kích thước đề xuất: 0.75\textwidth
\begin{figure}[H]
  \centering
  \fbox{\begin{minipage}{0.75\textwidth}\centering\vspace{2.5cm}%
    \textit{[Placeholder: Biểu đồ bar chart so sánh explored nodes]}%
  \vspace{2.5cm}\end{minipage}}
  \caption{Biểu đồ so sánh số node explored bởi 4 thuật toán (Greedy=15, A*=64, UCS=170, Dijkstra=$\geq 170$)}
  \label{fig:benchmark-chart}
\end{figure}

\section{Đánh giá UX/UI}

\subsection{Chế độ chọn ga}

Ưu điểm:
\begin{itemize}
  \item Dropdown danh sách 411 ga, sorted alphabet, dễ tìm.
  \item Hỗ trợ typeahead/search (gõ chữ cái đầu nhảy nhanh).
  \item Checkbox cấm tuyến/đoạn trực quan, không cần đọc hướng dẫn.
\end{itemize}

Hạn chế:
\begin{itemize}
  \item Dropdown rất dài (411 item), người dùng phải scroll, có thể khó sử dụng trên mobile.
  \item Một số ga tên chỉ là số (ví dụ \textquotedbl 1\textquotedbl, \textquotedbl 2\textquotedbl) do data OSM thiếu tag \texttt{name}.
\end{itemize}

\subsection{Chế độ click map}

Ưu điểm:
\begin{itemize}
  \item Trực quan, người dùng không cần nhớ tên ga.
  \item Cursor đổi thành crosshair, rõ ràng là chế độ pick.
  \item Tự động thoát mode sau khi pick.
\end{itemize}

Hạn chế:
\begin{itemize}
  \item Tìm ga gần nhất dùng Haversine thẳng, không tính đường bộ thực tế.
  \item Ngưỡng 50 km có thể quá lớn/nhỏ tuỳ khu vực.
\end{itemize}

\subsection{Bộ lọc hiển thị}

\begin{itemize}
  \item Toggle popover bằng nút "⚙ Hiển thị" ở góc trên-phải map.
  \item 20 checkbox tuyến, hỗ trợ master checkbox "Tất cả tuyến".
  \item Visual only, không ảnh hưởng pathfinding (để tránh confusion).
\end{itemize}

\subsection{Responsive Design}

\begin{itemize}
  \item Desktop (1280+): Sidebar form bên trái, map toàn màn hình bên phải.
  \item Tablet (768-1279): Stacked layout, form trên, map dưới.
  \item Mobile (360-767): Fullscreen form/map, collapse section.
  \item Touch-friendly: Button min-height 40px, 44px primary button.
\end{itemize}

\subsection{Accessibility}

\begin{itemize}
  \item ARIA labels trên button (aria-pressed, aria-expanded, aria-controls).
  \item Focus ring rõ ràng (focus-visible, custom --shadow-focus).
  \item Contrast AA trên body text (#14181c trên #ffffff ~ 17:1).
  \item prefers-reduced-motion support.
  \item Semantic HTML (button, form, role="dialog" trên popover).
\end{itemize}

\section{Performance Metrics}

\begin{itemize}
  \item \textbf{API response:} $< 20$ ms cho query điển hình (Hongqiao → Century Avenue).
  \item \textbf{Test suite:} 29 tests, $\approx 0.6$ s.
  \item \textbf{Initial load:} HTML + JS + CSS $< 100$ KB (gzip).
  \item \textbf{Data load:} JSON (206 KB) parse vào memory $< 50$ ms.
  \item \textbf{Polyline rendering:} Cache enabled, toggle blocked-line setStyle() không rebuild $\approx 10$ × tốc độ.
\end{itemize}

\section{Kết luận chương}

Chương này đã mô tả chi tiết cài đặt ứng dụng, từ công nghệ sử dụng, cấu trúc code, thực hiện 4 thuật toán, kết quả kiểm thử, benchmark, đến đánh giá UX/UI. Kết quả benchmark cho thấy A* là sự cân bằng tốt giữa tốc độ (64 node explored) và tính tối ưu. Greedy nhanh nhất (15 node) nhưng không tối ưu. UCS/Dijkstra đảm bảo tối ưu nhưng không dùng heuristic nên explored nhiều node hơn A*.
